

\begin{definition}\label{int_cv} \textbf{Integral de l\'inea para campos vectoriales.}
  Sea $C^{+} \subseteq  \Rn{n}$  una curva simple,  suave a trozos  (definición ~\ref{param}) orientada (definición ~ \ref{orien})  y sea $\boldsymbol{\sigma}:I\subseteq\R\to\mathbb{R}^n $ una parametrizaci\'on regular de $C$ que preserva su orientaci\'on  (definición ~\ref{orien_pre}). Sea $\mathbf{F}:A\subseteq\mathbb{R}^n \to\mathbb{R}^n $ un campo vectorial continuo sobre $C  \subseteq  A.$  Se define la \emph{integral de} $\mathbf{F}$ \emph{sobre la curva} $\boldsymbol{C}^{+}$ como
    \[
        \int_{C^{+}}\mathbf{F}\cdot d\mathbf{s}=\int_{\boldsymbol{\sigma}}\mathbf{F}\cdot d\mathbf{s} 
    \]
\end{definition}

\begin{obs} Bajo las condiciones de \ref{int_cv}, si $\boldsymbol{\beta}$ es una parametrizaci\'on regular de $C$ que no  preserva (o invierte) su orientaci\'on entonces: 
 \[
        \int_{C^{+}}\mathbf{F}\cdot d\mathbf{s}= - \int_{\boldsymbol{\beta}}\mathbf{F}\cdot d\mathbf{s} 
    \]
\end{obs}


